\section*{Capítulo \ref{ch:b2:wilson-cycle} --- Maciços antigos e o ciclo de Wilson}

\begin{solution}{exo:b2:wilson-cycle:1}
(a) $600 - 400 = \qty{200}{Myr}$, do rifteamento das margens até a última colisão.
(b) $480 - 400 = \qty{80}{Myr}$. (c) A orogenia caledoniana é um conjunto de fases distribuídas por
cerca de 200 milhões de anos, e cada fase começou em momentos diferentes em lugares diferentes: a palavra
localiza um oceano e seu cinturão, não um momento. Uma idade é dada pela escala estratigráfica ou em
milhões de anos.
\end{solution}

\begin{solution}{exo:b2:wilson-cycle:2}
(c) o rifte do leste da África, rifteamento; (a) o mar Vermelho, um oceano jovem; (d) o Atlântico, um oceano
largo com margens passivas; (e) a margem pacífica da América do Sul, subducção e
fechamento de um oceano; (b) o Himalaia, colisão; (f) a Bretanha, um cinturão aplainado, depois do qual
o rifteamento pode recomeçar.
\end{solution}

\begin{solution}{exo:b2:wilson-cycle:3}
(a) Depois da deposição das camadas do início do Siluriano, portanto depois de \qty{443.1}{Ma}, e antes da
deposição das camadas do Devoniano Superior, portanto antes de \qty{358.9}{Ma}, o fim do Devoniano.
(b) Menos de $443.1 - 358.9 = \qty{84.2}{Myr}$; na verdade menos, pois as camadas inferiores são mais jovens
que a base do Siluriano e as superiores mais antigas que o seu topo. (c) Deposição dos
turbiditos cinzentos no oceano; dobramento e clivagem nas colisões caledonianas; soerguimento e
erosão até uma superfície de terreno irregular; deposição das brechas e conglomerados vermelhos; suave
basculamento na orogenia varisca; a erosão atual, que corta a falésia.
\end{solution}

\begin{solution}{exo:b2:wilson-cycle:4}
(a) Os dois: um bloco de crosta continental separado dos continentes por oceanos enquanto
atravessava o Iapetus e, uma vez acoplado, um bloco limitado por falhas cuja história difere da de seus
vizinhos. (b) Nenhum: um continente, não um bloco acrescentado a um. (c) Um \omterm{def:b2:wilson-cycle:terrane}{terreno}, não um
\omterm{def:b2:wilson-cycle:terrane}{microcontinente}: sua crosta é a de um arco, não um fragmento de continente. (d) Nenhum: é
o limite entre \omterm{def:b2:wilson-cycle:terrane}{terrenos}, o traço do oceano fechado.
\end{solution}

\begin{solution}{exo:b2:wilson-cycle:5}
(a) $3000 - 2600 = 400$, $2600 - 1500 = 1100$, $1500 - 875 = 625$, $875 - 600 = 275$, $600 -
262.5 = 337.5$; média $2737.5/5 = \qty{547.5}{Myr}$. (b) $1100/275 = 4.0$. (c) $(1500 -
262.5)/3 = \qty{412.5}{Myr}$. A média muda de um quarto quando os pontos duvidosos são descartados,
e os espaçamentos variam por um fator de quatro: uma recorrência grosseira, não um período.
\end{solution}

\begin{solution}{exo:b2:wilson-cycle:6}
(a) $3000/(70 - 25) = \qty{66.7}{Myr}$. (b) $3000/70 = \qty{42.9}{Myr}$. (c) Depois que a dorsal
foi subductada, o fundo oceânico restante faz parte da placa distante, que carrega o continente
distante: a taxa de fechamento é a velocidade dessa placa em relação à fossa.
\end{solution}

\begin{solution}{exo:b2:wilson-cycle:7}
(a) $700/8 = \qty{87.5}{km}$; $700/30 = \qty{23.3}{km}$. (b) O cinturão varisco: a
\qty{30}{\celsius} por quilômetro, a crosta média funde; a \qty{8}{\celsius} por quilômetro,
a temperatura de fusão fica abaixo da base até mesmo de uma crosta espessada. (c) O gradiente baixo
descreve rochas levadas rapidamente para baixo numa zona de subducção, antes que pudessem se aquecer. Rochas que
ficam em profundidade numa crosta espessada se aquecem em direção a uma geoterma normal ao longo de milhões de anos, e podem
então atingir a fusão a uma profundidade menor.
\end{solution}

\begin{solution}{exo:b2:wilson-cycle:8}
(a) $\tau\ln(3/0.3) = 26.3\times2.303 = \qty{60.6}{Myr}$. (b) $290 - 60.6 = \qty{229}{Ma}$, no
Triássico Superior. (c) $\tau = (290 - 251.9)/\ln10 = \qty{16.5}{Myr}$. (d) O colapso
extensional da crosta espessada, que a afina e baixa sua superfície diretamente; também uma
taxa de erosão maior, um clima mais úmido ou rochas mais brandas, ou uma altitude inicial menor.
\end{solution}

\begin{solution}{exo:b2:wilson-cycle:9}
(a) A sutura, o traço do Iapetus, fica em terra na Grã-Bretanha e na Irlanda, a leste do
Atlântico; crosta da Laurentia, do lado oposto do Iapetus, fica perto de Rockall, do lado
europeu. O novo rifte cortou a oeste da antiga sutura: ele se abriu perto dela, não ao longo dela.
(b) Ao norte da sutura, embasamento e cobertura da margem da Laurentia; ao sul, os da
Avalônia, com histórias diferentes e, no início do Paleozoico, faunas diferentes. (c) O rifteamento
veio depois de 80 milhões de anos, mais do que os 30 a 50 depois dos quais uma zona resfriada pode ter se tornado
forte; no entanto, o Atlântico Central seguiu de fato os Apalaches. O calor de um cinturão se apaga; suas
falhas, tramas e composições permanecem e ainda podem guiar um rifte.
\end{solution}

\begin{solution}{exo:b2:wilson-cycle:10}
(a) $70 - 40/2 = \qty{50}{mm}$ por ano. (b) $\qty{1500}{km}/\qty{50}{km.Myr^{-1}} =
\qty{30}{Myr}$. (c) O oceano se estreita a $70 - 40 = \qty{30}{mm}$ por ano: $3000 - 30\times30 =
\qty{2100}{km}$. (d) Litosfera jovem e quente chega então à fossa, e nenhuma laje velha e fria
se desidrata mais sob o arco: o magmatismo de arco muda ou cessa.
\end{solution}

\begin{solution}{exo:b2:wilson-cycle:11}
(a) $h = \Delta T(\rho_m - \rho_c)/\rho_m = 25\times600/3300 = \qty{4.5}{km}$. (b) Todos os
\qty{25}{km} extras, 5.5 vezes a altura da cadeia. (c) $2700\times9.8\times25\,000 =
\qty{6.6e8}{Pa}$, cerca de \qty{0.66}{GPa}. (d) A cerca de \qty{0.7}{GPa} e às temperaturas da
raiz de uma crosta espessada, várias centenas de graus, a fácies anfibolito ou a fácies
granulito: os gnaisses e \omterm{def:b2:crustal-melting:migmatite}{migmatitos} dos maciços antigos.
\end{solution}

\begin{solution}{exo:b2:wilson-cycle:12}
(a) $S = 30\times10 = \qty{300}{km}$; $W_0 = 300 + 300 = \qty{600}{km}$; $T = 35\times600/300 =
\qty{70}{km}$. (b) $h = 35\times600/3300 = \qty{6.36}{km}$; $r = 4.5h = \qty{28.6}{km}$.
(c) $\Delta F = 2700\times9.8\times6364\times(35\,000 + 3182 + 14\,318) =
\qty{8.8e12}{N.m^{-1}}$, 1.4 vez o do Tibete. (d) Um cinturão assim é mais pesado que o Tibete, que já
se espalha: \omterm{def:b2:himalaya-andes:collapse}{colapso gravitacional}, por extensão de suas regiões altas e cavalgamento em suas bordas.
\end{solution}

\begin{solution}{pb:b2:wilson-cycle:1}
\textbf{1.} $720/(540 - 500) = \qty{18}{km}$ por milhão de anos, \qty{18}{mm} por ano.
\textbf{2.} \qty{36}{mm} por ano: uma dorsal de expansão lenta (20 a \qty{55}{mm} por ano).
\textbf{3.} $36\times(580 - 540) = \qty{1440}{km}$.
\textbf{4.} $36\times(580 - 500) = \qty{2880}{km}$.
\textbf{5.} Expansão simétrica a taxa constante a partir de largura zero a \qty{580}{Ma}, e nenhuma
subducção antes de \qty{500}{Ma}.
\textbf{6.} \qty{80}{Myr}, nas duas bordas do oceano, junto às margens continentais. É a
litosfera oceânica mais fria, mais espessa e mais densa do oceano, a mais pronta para afundar.
\textbf{7.} $64 - 36 = \qty{28}{mm}$ por ano.
\textbf{8.} $2880/28 = \qty{102.9}{Myr}$ depois de \qty{500}{Ma}: a \qty{397.1}{Ma}.
\textbf{9.} A dorsal, a \qty{1440}{km} da fossa, se aproxima dela a $64 - 18 =
\qty{46}{mm}$ por ano: depois de $1440/46 = \qty{31.3}{Myr}$, a \qty{468.7}{Ma}.
\textbf{10.} $2880 - 28\times31.3 = \qty{2004}{km}$.
\textbf{11.} $2004/40 = \qty{50.1}{Myr}$ depois: a $468.7 - 50.1 = \qty{418.6}{Ma}$.
\textbf{12.} Dentro dos 425 a \qty{400}{Ma} da colisão escandiana.
\textbf{13.} $40\times8 = \qty{320}{km}$.
\textbf{14.} $W_0 = 400 + 320 = \qty{720}{km}$; $T = 35\times720/400 = \qty{63.0}{km}$.
\textbf{15.} $h = 28\times600/3300 = \qty{5.09}{km}$; $r = 4.5\times5.09 = \qty{22.9}{km}$.
\textbf{16.} $\Delta F = 2700\times9.8\times5091\times(35\,000 + 2545 + 11\,455) =
\qty{6.60e12}{N.m^{-1}}$, cerca de \qty{8}{\percent} a mais que o do Tibete.
\textbf{17.} Parar de subir e se espalhar: extensão e falhas normais em seu alto interior,
cavalgamentos em suas bordas, como o Tibete faz hoje.
\textbf{18.} $f = 600/3300 = 0.182$; $k = 2.303/11 = 0.209$ por milhão de anos; $\tau = 1/(kf)
= \qty{26.3}{Myr}$.
\textbf{19.} $3.0\,e^{-30/26.3} = \qty{0.96}{km}$.
\textbf{20.} $26.3\ln(3.0/0.5) = \qty{47.1}{Myr}$ depois de \qty{410}{Ma}: a \qty{362.9}{Ma}.
\textbf{21.} A crosta extra é $h/f$: $3.0/0.182 = \qty{16.5}{km}$ cai para $0.5/0.182 =
\qty{2.75}{km}$, de modo que \qty{13.75}{km} de rocha são removidos.
\textbf{22.} $2700\times9.8\times13\,750 = \qty{3.6e8}{Pa}$, cerca de \qty{0.36}{GPa}.
\textbf{23.} Sim: no Devoniano Superior, o cinturão-modelo já desceu a cerca de \qty{1}{km} de altitude
média e continua caindo, com relevo restante entre seus vales e suas cristas; cascalhos vermelhos podem
se acumular em suas depressões bem antes de \qty{358.9}{Ma}, como em Siccar Point, onde repousam sobre uma superfície
com relevo.
\textbf{24.} \textbf{\qty[text-series-to-math=true]{2880}{km}}, a \qty{500}{Ma}, quando
a subducção começou.
\end{solution}
